\clearpage
\pdfbookmark[0]{Theorie der Kombination von Beobachtungen. Zweiter Teil}{theoria-posterior}
\thispagestyle{empty}
\SourceLeaf{032}
\begin{center}
\vspace*{2em}
{\Large THEORIE\par}
\bigskip
{\Large\bfseries DER KOMBINATION VON BEOBACHTUNGEN\par}
\bigskip
{\large DIE DEN KLEINSTEN FEHLERN AUSGESETZT IST\par}
\bigskip
{\large\textsc{Zweiter Teil}\par}
\bigskip
VON\par
\bigskip
{\large CARL FRIEDRICH GAUSS\par}
\bigskip
{\small DER KÖNIGLICHEN GESELLSCHAFT DER WISSENSCHAFTEN IM JAHR 1823, AM 2. FEBRUAR VORGELEGT.\par}
\bigskip
\rule{.22\linewidth}{.3pt}\par
\bigskip
Commentationes societatis regiae scientiarum Gottingensis recentiores.
\mbox{Vol.~\textsc{v}.}
\par\medskip
Göttingen MDCCCXXIII.
\par\bigskip
\rule{.22\linewidth}{.3pt}
\end{center}
\clearpage
\renewcommand{\Article}[1]{\par\medskip\Needspace{4\baselineskip}\pdfbookmark[1]{Art. #1}{theoria-posterior-art-#1}\centerline{\textbf{#1.}}\nobreak\smallskip}
\SourceLeaf{034}
\begin{center}
{\large THEORIE\par}
\medskip
{\Large\bfseries DER KOMBINATION VON BEOBACHTUNGEN\par}
\medskip
{\large DIE DEN KLEINSTEN FEHLERN AUSGESETZT IST.\par}
\medskip
{\large\textsc{Zweiter Teil.}\par}
\medskip
\rule{.22\linewidth}{.3pt}
\end{center}
\Article{23}
Es bleiben noch mehrere Untersuchungen, durch welche die vorangehende Theorie sowohl erläutert als auch erweitert werden wird.

Vor allem müssen wir untersuchen, ob die Elimination, durch die die Unbestimmten \(x,y,z\) etc. durch \(\xi,\eta,\zeta\) etc. ausgedrückt werden sollen, immer möglich ist. Da die Anzahl der ersteren derjenigen der letzteren gleich ist, ist aus der Theorie der Elimination bei linearen Gleichungen Folgendes über diese Elimination bekannt. Sind \(\xi,\eta,\zeta\) etc. voneinander unabhängig, so wird sie gewiss möglich sein; andernfalls unmöglich. Nehmen wir vorübergehend an, dass \(\xi,\eta,\zeta\) etc. nicht voneinander unabhängig sind. Vielmehr bestehe zwischen ihnen eine Identität
\[
0=F\xi+G\eta+H\zeta+\text{etc.}+K
\]
Wir werden also
\[
\begin{aligned}
F\Sigma aa+G\Sigma ab+H\Sigma ac+\text{etc.}&=0\\
F\Sigma ab+G\Sigma bb+H\Sigma bc+\text{etc.}&=0\\
F\Sigma ac+G\Sigma bc+H\Sigma cc+\text{etc.}&=0
\end{aligned}
\]
etc. haben sowie
\[
F\Sigma al+G\Sigma bl+H\Sigma cl+\text{etc.}=-K
\]
\par\SourceLeaf{035}
Setzen wir ferner
\[
\left.\begin{aligned}
aF+bG+cH+\text{etc.}&=\theta\\
a'F+b'G+c'H+\text{etc.}&=\theta'\\
a''F+b''G+c''H+\text{etc.}&=\theta''
\end{aligned}\right\}\quad\text{(I)}
\]
etc., so ergibt sich
\[
\begin{aligned}
a\theta+a'\theta'+a''\theta''+\text{etc.}&=0\\
b\theta+b'\theta'+b''\theta''+\text{etc.}&=0\\
c\theta+c'\theta'+c''\theta''+\text{etc.}&=0
\end{aligned}
\]
etc. sowie
\[
l\theta+l'\theta'+l''\theta''+\text{etc.}=-K
\]
Multiplizieren wir also die Gleichungen (I) jeweils mit \(\theta,\theta',\theta''\) etc. und addieren, so erhalten wir:
\[
0=\theta\theta+\theta'\theta'+\theta''\theta''+\text{etc.}
\]
Diese Gleichung kann offenbar nur bestehen, wenn zugleich \(\theta=0,\allowbreak\theta'=0,\allowbreak\theta''=0\) etc. gilt. Hieraus schließen wir zunächst, dass notwendig \(K=0\) sein muss. Sodann zeigen die Gleichungen (I), dass die Funktionen \(v,v',v''\) etc. so beschaffen sind, dass ihre Werte unverändert bleiben, wenn die Werte der Größen \(x,y,z\) etc. Zuwächse oder Abnahmen erhalten, die jeweils \(F,G,H\) etc. proportional sind; offenbar gilt dasselbe für die Funktionen \(V,V',V''\) etc. Die Annahme kann also nur in einem Fall bestehen, in dem es selbst aus den genauen Werten der Größen \(V,V',V''\) etc. unmöglich gewesen wäre, die Werte der Unbekannten \(x,y,z\) etc. zu bestimmen, also in dem die Aufgabe ihrer Natur nach unbestimmt gewesen wäre. Diesen Fall haben wir von unserer Untersuchung ausgeschlossen.
\Article{24}
Mit \(\beta,\beta',\beta''\) etc. bezeichnen wir Multiplikatoren, die zur Unbestimmten \(y\) dieselbe Beziehung haben wie \(\alpha,\alpha',\alpha''\) etc. zu \(x\); es sei nämlich
\[
\begin{aligned}
a[\beta\alpha]+b[\beta\beta]+c[\beta\gamma]+\text{etc.}&=\beta\\
a'[\beta\alpha]+b'[\beta\beta]+c'[\beta\gamma]+\text{etc.}&=\beta'\\
a''[\beta\alpha]+b''[\beta\beta]+c''[\beta\gamma]+\text{etc.}&=\beta''
\end{aligned}
\]
etc., sodass identisch gilt
\[
\beta v+\beta'v'+\beta''v''+\text{etc.}=y-B
\]
\par\SourceLeaf{036}
Ebenso seien \(\gamma,\gamma',\gamma''\) etc. entsprechende Multiplikatoren bezüglich der Unbestimmten \(z\), nämlich
\[
\begin{aligned}
a[\gamma\alpha]+b[\gamma\beta]+c[\gamma\gamma]+\text{etc.}&=\gamma\\
a'[\gamma\alpha]+b'[\gamma\beta]+c'[\gamma\gamma]+\text{etc.}&=\gamma'\\
a''[\gamma\alpha]+b''[\gamma\beta]+c''[\gamma\gamma]+\text{etc.}&=\gamma''
\end{aligned}
\]
etc., sodass identisch
\[
\gamma v+\gamma'v'+\gamma''v''+\text{etc.}=z-C
\]
und so weiter gilt. \Needspace{16\baselineskip}Auf diese Weise werden wir, ebenso wie wir bereits in Art. 20
\[
\Sigma\alpha a=1,\quad\Sigma\alpha b=0,\quad\Sigma\alpha c=0,\quad\text{etc., sowie}\quad\Sigma\alpha l=-A
\]
gefunden hatten, auch
\[
\begin{aligned}
\Sigma\beta a&=0,\quad\Sigma\beta b=1,\quad\Sigma\beta c=0,\quad\text{etc., und}\quad\Sigma\beta l=-B\\
\Sigma\gamma a&=0,\quad\Sigma\gamma b=0,\quad\Sigma\gamma c=1,\quad\text{etc., und}\quad\Sigma\gamma l=-C
\end{aligned}
\]
und so weiter haben. Ebenso wie sich in Art. 20 \(\Sigma\alpha\alpha=[\alpha\alpha]\) ergab, wird ferner
\[
\Sigma\beta\beta=[\beta\beta],\quad\Sigma\gamma\gamma=[\gamma\gamma]\quad\text{etc.}
\]
sein. Multiplizieren wir nun die Werte von \(\alpha,\alpha',\alpha''\) etc. (Art. 20. IV) jeweils mit \(\beta,\beta',\beta''\) etc. und addieren, so erhalten wir
\[
\alpha\beta+\alpha'\beta'+\alpha''\beta''\quad\text{etc.}=[\alpha\beta],\quad\text{oder}\quad\Sigma\alpha\beta=[\alpha\beta]
\]
Multiplizieren wir dagegen die Werte von \(\beta,\beta',\beta''\) etc. jeweils mit \(\alpha,\alpha',\alpha''\) etc. und addieren, so ergibt sich ebenso
\[
\alpha\beta+\alpha'\beta'+\alpha''\beta''+\text{etc.}=[\beta\alpha],\quad\text{also}\quad[\alpha\beta]=[\beta\alpha]
\]
Ganz genauso ergibt sich
\[
[\alpha\gamma]=[\gamma\alpha]=\Sigma\alpha\gamma,\quad[\beta\gamma]=[\gamma\beta]=\Sigma\beta\gamma\quad\text{etc.}
\]
\clearpage
